% =========================================================================
% Capítulo 1: Cinemática
% Tema 02: Coordenadas esféricas
% Problema 04: Anillo en barra giratoria
% Origen: Ayudantía 2, 2026
% =========================================================================
\subsubsection{Anillo en barra giratoria}
Un anillo se desliza a lo largo de una barra, como se muestra en la Figura \ref{fig:c1_e02_p04_barra_giratoria}. La barra gira libremente alrededor de una bisagra en un punto fijo $O$. Simultáneamente, la barra gira alrededor del eje vertical con una velocidad angular constante $\vec{\omega}_1$ respecto al suelo. Sabiendo que $r$ describe la ubicación del anillo $P$ a lo largo de la varilla, $\phi$ es el ángulo medido desde la vertical y $\|\vec{\omega}_1\| = \omega_1$, determine:

\begin{enumerate}
    \item Los marcos de referencia y los sistemas de coordenadas.
    \item La velocidad angular del anillo según un observador fijo en el suelo.
    \item La aceleración angular del suelo según la ve un observador fijo en el anillo.
    \item La velocidad del anillo según un observador fijo en el suelo.
    \item La aceleración del anillo según un observador fijo en el suelo.
\end{enumerate}

\begin{figure}[!ht]
    \centering
    \begin{tikzpicture}[>=Latex, scale=1.5]
        % --- PARÁMETROS ---
        \def\Lb{3.0}     % longitud de la barra [u]
        \def\phiAng{50}  % ángulo phi medido desde la vertical [°]

        % --- SOPORTES Y ELEMENTOS FIJOS ---
        \fill[pattern=north east lines] (-0.3,-0.9) rectangle (0.3,-0.2);
        \draw[thick] (-0.3,-0.2) -- (-0.3,-0.9) -- (0.3,-0.9) -- (0.3,-0.2);
        \draw[thick] (-0.3,-0.2) -- (0.3,-0.2);

        % --- GEOMETRÍA PRINCIPAL ---
        \coordinate (O) at (0,0);
        \coordinate (P) at ({90-\phiAng}:\Lb);
        \begin{scope}[rotate={90-\phiAng}]
            \shadedraw[left color=gray!10, right color=gray!40, draw=black, rounded corners=1.5pt]
                (-0.1,-0.06) rectangle ({\Lb+0.3},0.06);
            % anillo y mango
            \draw[thick, fill=gray!30] ({\Lb-0.12},-0.09) rectangle ({\Lb+0.12},0.09);
        \end{scope}

        % --- SISTEMAS DE COORDENADAS Y REFERENCIAS ---
        \draw[dotted] (O) -- (-0.9,0);
        \draw[dotted] (O) -- (3.3,0);
        \draw[dashed] (O) -- (0,2.6);

        % --- PUNTOS Y COORDENADAS ---
        \filldraw[fill=gray!80] (O) circle (2pt) node[below left] {$O$};
        \filldraw (P) circle (1.3pt) node[above left] {$P$};

        % --- VECTORES Y ÁNGULOS ---
        \draw[->, very thick] (-0.9,0.3) -- (-0.9,2.6) node[above] {$\omega_1$};
        \draw[->, thick] (-1.15,1.9) arc (190:-10:0.25 and 0.09);
        \draw[->] (0,1.0) arc (90:{90-\phiAng}:1.0);
        \node at ({90-0.5*\phiAng}:1.25) {$\phi$};

        % --- ANOTACIONES ---
        \node[above left] at ({90-\phiAng}:{0.55*\Lb}) {$r$};
        \draw[<->, thin] ($({90-\phiAng}:0.0)+({90-\phiAng-90}:-0.2)$) -- ($(P)+({90-\phiAng-90}:-0.2)$);
    \end{tikzpicture}
    \caption{Anillo $P$ deslizando sobre una barra que gira.}
    \label{fig:c1_e02_p04_barra_giratoria}
\end{figure}

% --- SOLUCIÓN ---------------------------------------------------------------
\begin{solucion}{Anillo en barra giratoria}
\textbf{Marcos de referencia y Sistemas de coordenadas}\\
Se definen tres marcos de referencia para aplicar el método de dos bases cilíndricas (evitando así las coordenadas esféricas directas):
\begin{itemize}
    \item $\mathcal{N}$ El suelo (inercial) $\longrightarrow\{\vec{E}_x;\vec{E}_y;\vec{E}_z\}$ donde $\vec{E}_z$ es el eje vertical de rotación.
    \item $\mathcal{A}$ Marco azimutal (gira con la barra) $\longrightarrow\{\vec{e}_R;\vec{e}_\theta;\vec{E}_z\}$
    \item $\mathcal{B}$ La barra $\longrightarrow\{\vec{e}_r;\vec{e}_\phi;\vec{e}_\theta\}$
\end{itemize}
El plano formado por $\vec{e}_R$ y $\vec{E}_z$ contiene siempre a la varilla. El marco $\mathcal{B}$ está rotado un ángulo $\phi$ respecto a $\vec{E}_z$ dentro de ese mismo plano vertical.
La relación geométrica para proyectar vectores desde $\mathcal{A}$ hacia $\mathcal{B}$ es:
$$\vec{E}_z = \cos(\phi)\vec{e}_r - \sin(\phi)\vec{e}_\phi$$
$$\vec{e}_R = \sin(\phi)\vec{e}_r + \cos(\phi)\vec{e}_\phi$$

\patotip[squareglasses]{Una base esférica no es más que \emph{dos} cilíndricas encadenadas: primero se gira $\omega_1$ en torno a $\vec{E}_z$ (marco $\mathcal{A}$) y luego $\dot\phi$ en torno a $\vec{e}_\theta$ (marco $\mathcal{B}$). Cada marco aporta una velocidad angular simple y después se suman.}

\textbf{Velocidad angular del anillo según un observador fijo en el suelo}\\
Por propiedad aditiva de velocidades angulares:
$$\velang{N}{B} = \velang{N}{A} + \velang{A}{B}$$
Sabiendo que el marco $\mathcal{A}$ gira a velocidad constante $\omega_1$ en $\vec{E}_z$ y la varilla ($\mathcal{B}$) cae con ángulo $\phi$ en la dirección $\vec{e}_\theta$:
$$\velang{N}{A} = \omega_1\vec{E}_z$$
$$\velang{A}{B} = \dot\phi\vec{e}_\theta$$
Proyectando el vector $\vec{E}_z$ a la base solidaria de la barra $\mathcal{B}$:
$$\velang{N}{B} = \omega_1(\cos(\phi)\vec{e}_r - \sin(\phi)\vec{e}_\phi) + \dot\phi\vec{e}_\theta = \omega_1\cos(\phi)\vec{e}_r - \omega_1\sin(\phi)\vec{e}_\phi + \dot\phi\vec{e}_\theta$$

\textbf{Aceleración angular del suelo según la ve un observador fijo en el anillo}\\
Nos piden $\accang{B}{N}$, y por teoría sabemos que $\accang{B}{N} = -\accang{N}{B}$.
Calculamos primero $\accang{N}{B}$ aplicando la derivada:
$$\accang{N}{B} = \der{}{N}\left(\velang{N}{B}\right) = \der{(\omega_1\vec{E}_z)}{N}{} + \der{(\dot\phi\vec{e}_\theta)}{N}{}$$
Como $\omega_1$ es constante y $\vec{E}_z$ es fijo en $\mathcal{N}$, el primer término es cero. Para el segundo, usamos Teorema del Transporte respecto a $\mathcal{A}$:
$$\der{(\dot\phi\vec{e}_\theta)}{N} = \transporta{(\dot\phi\vec{e}_\theta)}{N}{A}$$
$$\der{(\dot\phi\vec{e}_\theta)}{N} = \ddot\phi\vec{e}_\theta + \velang{N}{A}\times(\dot\phi\vec{e}_\theta) = \ddot\phi\vec{e}_\theta + (\omega_1\vec{E}_z)\times(\dot\phi\vec{e}_\theta) = \ddot\phi\vec{e}_\theta - \omega_1\dot\phi\vec{e}_R$$
Proyectamos $\vec{e}_R$ en la base $\mathcal{B}$:
$$\accang{N}{B} = \ddot\phi\vec{e}_\theta - \omega_1\dot\phi(\sin(\phi)\vec{e}_r + \cos(\phi)\vec{e}_\phi)$$
Invirtiendo el signo obtenemos lo solicitado:
\respuesta{$\displaystyle\accang{B}{N} = \omega_1\dot\phi\sin(\phi)\vec{e}_r + \omega_1\dot\phi\cos(\phi)\vec{e}_\phi - \ddot\phi\vec{e}_\theta$}

\textbf{Velocidad del anillo según un observador fijo en el suelo}\\
El vector posición del anillo es simplemente $\vec{r}_{P/O} = r\vec{e}_r$. Aplicamos el Teorema del Transporte respecto al marco $\mathcal{B}$:
$$\marc{\vec{v}_P}{N} = \der{(r\vec{e}_r)}{N} = \transporta{(r\vec{e}_r)}{N}{B}$$
$$\marc{\vec{v}_P}{N} = \dot{r}\vec{e}_r + \velang{N}{B}\times(r\vec{e}_r)$$
Calculamos el producto cruz:
$$\velang{N}{B}\times(r\vec{e}_r) = \left(\omega_1\cos(\phi)\vec{e}_r - \omega_1\sin(\phi)\vec{e}_\phi + \dot\phi\vec{e}_\theta\right) \times (r\vec{e}_r) = \omega_1 r \sin(\phi)\vec{e}_\theta + r\dot\phi\vec{e}_\phi$$
Sumando la derivada local:
\respuesta{$\displaystyle\marc{\vec{v}_P}{N} = \dot{r}\vec{e}_r + r\dot\phi\vec{e}_\phi + r\omega_1\sin(\phi)\vec{e}_\theta$}

\textbf{Aceleración del anillo según un observador fijo en el suelo}\\
Volvemos a aplicar el Teorema del Transporte, esta vez sobre la velocidad, apoyándonos en el marco $\mathcal{B}$:
$$\marc{\vec{a}_P}{N} = \der{\vec{v}_P}{N} = \transporta{\vec{v}_P}{N}{B}$$

Paso 1: Derivada local de la velocidad en $\mathcal{B}$:
$$\der{\vec{v}_P}{B} = \ddot{r}\vec{e}_r + (\dot{r}\dot\phi + r\ddot\phi)\vec{e}_\phi + (\dot{r}\omega_1\sin(\phi) + r\omega_1\dot\phi\cos(\phi))\vec{e}_\theta$$

Paso 2: Producto cruz de arrastre ($\velang{N}{B}\times\vec{v}_P$):
$$ \velang{N}{B}\times\vec{v}_P = \begin{vmatrix} \vec{e}_r & \vec{e}_\phi & \vec{e}_\theta \\ \omega_1\cos(\phi) & -\omega_1\sin(\phi) & \dot\phi \\ \dot{r} & r\dot\phi & r\omega_1\sin(\phi) \end{vmatrix} $$

Resolviendo por componentes:
\begin{itemize}
    \item En $\vec{e}_r$: $-\omega_1^2 r \sin^2(\phi) - r\dot\phi^2$
    \item En $\vec{e}_\phi$: $\dot{r}\dot\phi - r\omega_1^2\sin(\phi)\cos(\phi)$
    \item En $\vec{e}_\theta$: $\omega_1 r \dot\phi \cos(\phi) + \dot{r}\omega_1\sin(\phi)$
\end{itemize}

Agrupando ambos pasos y factorizando, se obtiene la aceleración total:
\respuesta{$\displaystyle\marc{\vec{a}_P}{N} = \left(\ddot{r} - r\dot\phi^2 - r\omega_1^2\sin^2(\phi)\right)\vec{e}_r + \left(2\dot{r}\dot\phi + r\ddot\phi - r\omega_1^2\sin(\phi)\cos(\phi)\right)\vec{e}_\phi + \left(2\dot{r}\omega_1\sin(\phi) + 2r\omega_1\dot\phi\cos(\phi)\right)\vec{e}_\theta$}
\end{solucion}
